\begin{problem}[Generic and closed points in product rings]
Let
\[
A=A_1\times\cdots\times A_r.
\]
Describe $\Spec A$ as a disjoint union of the spectra of the factors.  Determine how closed
points, generic points of irreducible components, and specialization relations are read
factor by factor.  Specialize to $A=k^r$.
\end{problem}

\paragraph{Why this problem matters.}
Products show how disconnected geometry interacts with generic points: there is generally
one generic point per irreducible component, not one for the whole space.

\paragraph{Useful ideas before starting.}
Prime ideals of a finite product come from exactly one factor, while the other factors are
entirely present in the ideal.

\paragraph{Guided hints.}
Use
\[
\Spec(A_1\times A_2)\cong\Spec A_1\sqcup\Spec A_2
\]
and iterate.

\begin{solution}
For a finite product,
\[
\Spec A
\cong
\bigsqcup_{i=1}^r\Spec A_i.
\]
Each summand is both open and closed.  Therefore all specialization relations remain
inside a single factor; points in distinct factors are incomparable.

A point is closed in $\Spec A$ exactly when it is closed in its factor, hence corresponds
to a maximal ideal of that factor.  Irreducible components of $\Spec A$ are precisely the
irreducible components of all the factor spectra, and their generic points are the
corresponding factorwise minimal primes.

If $A=k^r$, each factor spectrum is a single point.  Hence $\Spec k^r$ is a discrete set of
$r$ points.  Every point is closed and is also the generic point of its singleton
irreducible component.  For $r>1$ there is no generic point of the whole space because it
is reducible.
\end{solution}

\paragraph{What happened mathematically.}
Disconnected product decomposition separates specialization orders completely.

\paragraph{Extensions.}
Relate the clopen components to the standard idempotents
\[
(1,0,\dots,0),\ldots,(0,\dots,0,1).
\]

\paragraph{Provenance.}
New pointwise synthesis of the product/idempotent legacy material routed primarily to
VI/05.
